% CS F407 Artificial Intelligence: final project report template.
% Limit: 2 pages, including everything except references.
% Do not change the font size, margins or spacing.
\documentclass[11pt,a4paper,twocolumn]{article}

\usepackage[margin=2cm,columnsep=0.8cm]{geometry}
\usepackage[T1]{fontenc}
\usepackage{lmodern}
\usepackage{microtype}
\usepackage{amsmath,amssymb}
\usepackage{graphicx}
\usepackage{booktabs}
\usepackage{tabularx}
\usepackage{enumitem}
\usepackage{titlesec}
\usepackage[hidelinks]{hyperref}

\titleformat{\section}{\normalsize\bfseries}{\thesection}{0.6em}{}
\titlespacing*{\section}{0pt}{1.2ex plus 0.3ex}{0.6ex}
\setlist{nosep,leftmargin=*}
\setlength{\parskip}{0.3ex}

\newcommand{\coursename}{CS F407 Artificial Intelligence, AY 2026--27 Semester I}
\newcommand{\teamno}{Team 404byArcticDonkeys}
\newcommand{\githuburl}{https://github.com/Naitik-Gupta96/AI_Project}

\begin{document}

\twocolumn[{%
\begin{center}
  {\large\bfseries Intelligible Blackboard Architecture with Counterfactual Agents}\\[1ex]
  \teamno\\
  Naitik Gupta (2024A7PS0603G), Malay Neema (2024A1PS0286G), Mayank Sharma (2024A7PS0627G),
  Yash Singh (2024A7PS0617G)\\[0.5ex]
  {\small \coursename, BITS Pilani, K.K. Birla Goa Campus}\\[1.5ex]
  \begin{minipage}{0.85\textwidth}
    \small
    \centerline{\bfseries Abstract}\vspace{0.5ex}
    Language-model agents that pass text along a chain lose context and compound each other's
    errors. We built a system in which three agents coordinate only through a shared blackboard,
    and every message follows the Two-Way Intelligibility Protocol (PXP). When the board
    deadlocks, counterfactual-capable agents replay the session in a sandbox with alternatives
    to their own past messages, credit each alternative by how much closer it brings the team to
    consensus, and apply the best one to the live board. The system comprises an event-sourced
    blackboard, an event-driven scheduler, Qwen2.5-7B-Instruct agents, a live web interface and
    loaders for MSCoRe, KramaBench and MedAgentBench. In a reduced-scale ablation (4 groups
    $\times$ 25 tasks), every session reached consensus without a deadlock, so all groups scored
    the same (convergence success 0.16). With scripted agents on synthetic tasks that do
    deadlock, convergence success rises from 0.490 to 0.544 as the share of counterfactual
    agents grows from 0 to 100\%.
  \end{minipage}
\end{center}
\vspace{1.5ex}
}]

\section{Introduction}
Multi-agent systems built from language models usually pass instructions along a chain. Over
long tasks such chains lose context, and an early mistake is passed on as fact. A blackboard
architecture removes the chain: agents read and write one shared, structured state and react to
its changes. A shared board alone, however, does not make the interaction intelligible, and
agents that disagree can stop making progress.

We address both problems. The \emph{input} is a task: a question with optional answer options,
public context, private evidence for each role, and tools (a Python sandbox over a data lake, or
a FHIR server of patient records). The \emph{output} is the team's agreed answer, the complete
log of PXP messages with its intelligibility class, and the token cost. We contribute:
\begin{itemize}
\item a thread-safe, event-sourced blackboard that enforces PXP on every write, with an
  event-driven scheduler;
\item retrospective counterfactual credit assignment that changes the live board through the
  agent's prompt and its voting weight;
\item a web interface that streams the board and its counterfactual replays;
\item an ablation over the share of counterfactual-capable agents on three public benchmarks.
\end{itemize}

\section{Background and related work}
\textbf{Blackboards.} Hearsay-II~\cite{erman1980} coordinated independent knowledge sources
through a shared data structure and a control component that decides who acts next.

\textbf{PXP.} The Two-Way Intelligibility Protocol~\cite{baskar2023} governs agents that
predict and explain. Mestha et al.~\cite{mestha2025} give it a multi-turn form (their
Algorithm~1). Each message carries a prediction, an explanation and a tag: \textsc{ratify},
\textsc{revise}, \textsc{refute} or \textsc{reject}. An interaction is classified as One-Way,
Strong or Ultra-Strong intelligible (their Definitions~1--2).

\textbf{LLM teams and counterfactual credit.} Multi-agent debate improves the reasoning of
language models~\cite{du2024debate}, but it has no shared state and no protocol that can be
checked. In cooperative multi-agent reinforcement learning, COMA~\cite{foerster2018coma} assigns
credit by comparing an agent's action with a baseline over its alternatives. We do not learn a
critic; we replay the recorded session from each alternative message.

\textbf{Benchmarks.} We use MSCoRe~\cite{mscore} (multi-stage reasoning, in Chinese),
KramaBench~\cite{kramabench} (data-to-insight pipelines over data lakes) and
MedAgentBench~\cite{medagentbench} (clinical tasks against a FHIR server).

\section{Method}
\textbf{PXP messages.} Message $i$ carries a prediction $y$ and an explanation $e$. Its sender
compares the peer message it answers, $(y',e')$, with its own previous position $(y'',e'')$:
\begin{gather*}
m=\mathrm{MATCH}(y',y''),\qquad a=\mathrm{AGREE}(e',e''),\\
c=\neg\mathrm{MATCH}(y,y'')\lor\neg\mathrm{AGREE}(e,e'').
\end{gather*}
The tag follows from these judgments:
\begin{itemize}
\item message~1 is \textsc{init};
\item $m\land a$ gives \textsc{ratify};
\item exactly one of $m$, $a$ gives \textsc{revise} if $c$, and \textsc{refute} otherwise;
\item neither gives \textsc{reject} when $i>k$ ($k=3$), and \textsc{revise}/\textsc{refute} before.
\end{itemize}
The model supplies the judgments that need understanding: AGREE, whether its reasoning changed,
and MATCH for free-form answers. Code computes MATCH over listed options, then the tag. A
grammar checks every message on commit and rejects (never repairs) an invalid one.

\textbf{Blackboard and scheduler.} Every change is an event, and the state is the fold of the
event log, so live writes, the step-back inspector and the counterfactual sandboxes see
identical states. Writes are serialized and versioned: a write based on a stale version raises a
conflict instead of overwriting. Logs persist as fsynced JSON lines or Redis streams. The
scheduler wakes on each state change, checks for consensus and deadlock, and ranks the eligible
agents: $+10$ for a pending intervention, $+3$ if the agent has not spoken, $+2$ for news since
its last post, $+1$ if its last tag was not \textsc{ratify}, with ties broken by registration
order so that replays are deterministic. Consensus holds when at least two agents ratify the
same anchor message and hold at least $\tau=0.75$ of the voting weight. A deadlock is a
\textsc{reject}, three \textsc{refute}s in a row, or six messages without a new prediction.

\textbf{Agents.} Each benchmark has a team of three personas: one expert per value-chain stage
for MSCoRe; a data engineer, an analyst and a verifier for KramaBench; a retriever, a clinician
and a safety reviewer for MedAgentBench. An agent sees the task, its private evidence, the
message it answers, its previous position, a clipped history, its lessons and its weight. The
model returns JSON with the fields \texttt{prediction}, \texttt{explanation},
\texttt{explanation\_agrees}, \texttt{changed} and, for free-form answers,
\texttt{prediction\_matches}, with decoding constrained to that schema. Before posting, an agent
may make up to four tool calls: \texttt{list\_files} and a sandboxed \texttt{python} on
KramaBench, \texttt{fhir\_get} on MedAgentBench.

\textbf{Counterfactual credit assignment.} On a deadlock, each capable agent reflects on its last
three messages and proposes, for each, up to three alternatives (another prediction, another
wording of its explanation, or both) and a lesson. Each alternative replaces the original
message in a sandbox rebuilt from the event log, and clones of all agents play the session to
the end. With proximity $P$ the weighted share of the largest group ratifying one thesis,
\[
\mathrm{credit}=P_{\text{sandbox end}}-P_{\text{live, at deadlock}}.
\]
The alternative with the highest positive credit is applied to the live board. Its lesson joins
the agent's system prompt, the agent's position is revised, and its weight becomes
$w\leftarrow\max(0.25,\,w(1-0.5\cdot\mathrm{credit}))$. The board then returns from
\textsc{deadlock} to \textsc{running}; at most two interventions are allowed per session.

For example, on a synthetic task the analyst opens with Option~A, which the critic and the
specialist each know to be infeasible, and the board deadlocks. In the sandbox, replacing the
opening with Option~B ends in consensus (credit $+0.67$). The analyst's weight drops to 0.67, and
the live session then agrees on B.

\textbf{Interface.} A FastAPI server streams each session over a WebSocket to a React/D3 view.
The task is the central node, and messages branch out from it by tag. Cascading
\textsc{reject}s flash red, a split panel shows the counterfactual timelines beside the live
one, and an inspector steps forward and back through any session.

\section{Experimental setup}
\textbf{Design.} The brief's ablation has four groups, in which 0\% (control), 33\%, 67\% and
100\% of the three agents are counterfactual-capable; capable agents rotate with the task index.
Every task runs in every group, so the groups are paired.

\textbf{Compute and scale.} The brief specifies 500 tasks per group, and we built the system to
serve Qwen2.5-7B-Instruct in bf16 through vLLM on a GPU. The bf16 weights alone take 15.2 GB,
which needs a 24 GB GPU. We could not get one, so we ran on the best machine we had, an Apple M5
laptop, with a 4-bit build of the same model. A 20-session smoke run there first exposed several
bugs in the benchmark adapters, which we fixed. At the laptop's speed the full design would take
about 12 hours. We therefore kept the design and reduced only its size: 25 tasks per group,
sampled with seed~0, 100 sessions in all. These are 10 MSCoRe tasks stratified by value chain
and difficulty (accurate when ROUGE-L F1~$\ge0.3$), 5 KramaBench tasks (the benchmark's
per-type metrics) and 10 MedAgentBench tasks, one per task type (the official grader). The
repository's full-scale configuration runs the specified design unchanged.

\textbf{Model and hardware.} We used Qwen2.5-7B-Instruct~\cite{qwen25} in a 4-bit build
(Q4\_K\_M) through Ollama 0.40.3, at temperature~0 and seed~0 with an 8,192-token context. A
vLLM~\cite{kwon2023vllm} path is also implemented. Sessions allow 12 messages and at most two
interventions. The run took 37 minutes on 24~Nov~2026 on an Apple M5 laptop with 16 GB, with
the MedAgentBench FHIR server in Docker.

\textbf{Metrics.} \emph{Convergence success} is consensus within the budget together with an
accurate answer, with 95\% Wilson intervals. We also report deadlock and recovery rates, and
\emph{intelligibility depth} (the share of Ultra-Strong, Strong, One-Way and None sessions).
\emph{Token cost} is the mean of prompt plus completion tokens per instance, on the live board
and in the sandbox. Each group is compared with the control by an exact McNemar test.

\textbf{Pipeline check.} To test the counterfactual layer where deadlocks do occur, we also ran
the full $4\times500$ design on synthetic decision tasks, in which each role holds private
evidence that rules options out. These agents are deterministic and rule-based, not language
models, so the numbers describe the mechanism, not LLM behaviour.

\section{Results}
\begin{table}[t]
\centering
\small
\begin{tabular}{@{}lrcr@{}}
\toprule
Benchmark & $n$ & Success [95\% CI] & Tokens \\
\midrule
MSCoRe & 10 & 0.00 [0.00, 0.28] & 3.6k \\
KramaBench & 5 & 0.20 [0.04, 0.62] & 9.0k \\
MedAgentBench & 10 & 0.30 [0.11, 0.60] & 14.1k \\
All & 25 & 0.16 [0.06, 0.35] & 8.9k \\
\bottomrule
\end{tabular}
\caption{Qwen2.5-7B run, 25 tasks per group. Convergence success (higher is better) is identical
in all four groups. Tokens: mean per instance in the control group (lower is better); the other groups differ by at most 2\%.}
\label{tab:results}
\end{table}

\begin{table}[t]
\centering
\small
\begin{tabular}{@{}lccccc@{}}
\toprule
Group & Dead. & Recov. & Conv. & Calls & $p$ \\
\midrule
0\% & 0.51 & 0.000 & 0.490 & 4.27 & -- \\
33\% & 0.51 & 0.024 & 0.502 & 6.67 & 0.031 \\
67\% & 0.51 & 0.075 & 0.528 & 9.10 & $4\cdot10^{-6}$ \\
100\% & 0.51 & 0.106 & 0.544 & 11.85 & $2\cdot10^{-8}$ \\
\bottomrule
\end{tabular}
\caption{Pipeline check with scripted agents, 500 synthetic tasks per group. Dead.: share of
sessions that deadlocked. Recov.: share of those that still reached consensus. Conv.:
convergence success (higher is better). Calls: model calls per session (lower is better).
$p$: McNemar test against the control.}
\label{tab:mock}
\end{table}

\textbf{Language-model run} (Table~\ref{tab:results}). Every session in every group reached
consensus in four messages: an \textsc{init}, then a \textsc{ratify} from each of the other two
agents. No session deadlocked, so no intervention ran. The four groups gave the same answer on
all 25 tasks, every McNemar test has $b=c=0$ ($p=1$), and sandbox tokens are zero. Every session
is Strong and none is Ultra-Strong, since no agent ever revised. MSCoRe answers of 49 to 91
characters against long reference essays give ROUGE-L precision 0.69 but recall 0.04 (F1 0.08),
so none passes the 0.3 threshold. KramaBench solved 1 of 5 tasks, and MedAgentBench 3 of 10
(task types 4, 5 and 9).

\textbf{Pipeline check} (Table~\ref{tab:mock}). Half of the synthetic sessions deadlock; in 134
of the 255 deadlocked tasks no option is feasible, so deadlock is the correct outcome. As the
share of capable agents grows, more deadlocks are recovered and convergence success rises from
0.490 to 0.544. No capable group did worse than the control on any task, and the McNemar tests
reject equality at every level. The replays are not free: model calls per session rise from
4.27 to 11.85.

\section{Discussion and limitations}
The language-model run does not answer the research question, because the situation that
counterfactual agents act on never arose. The 7B model ratified the opening message in every
session, even when that answer was wrong, so the teams agreed too easily rather than
deadlocking. The personas gave the agents different viewpoints but no conflicting information,
whereas in the synthetic tasks each role privately knows something the others do not, and
deadlocks follow. The pipeline check shows that, when deadlocks occur, the mechanism works and
helps. It also shows that which agent is capable matters: a deadlock is resolved only when the
agent whose message caused it can reconsider that message. Replays add calls and tokens, so on
this evidence counterfactual reasoning buys convergence at a cost rather than saving tokens.

\textbf{Limitations.} With 25 tasks per group the intervals are up to 58 points wide, and one
4-bit 7B model ran on a laptop. MSCoRe's criterion penalizes our short answers, and KramaBench's
paraphrase judge is the same model as the agents.

\textbf{Next steps.} Run the full design on a GPU (the configuration is ready), build benchmark
variants that partition evidence and tools between roles so that agents genuinely disagree, and
test larger models.

\section{Conclusion}
We built every part of the brief: a PXP-governed blackboard, an event-driven scheduler,
counterfactual credit assignment applied to the live board, a live interface and three benchmark
pipelines. On real benchmarks a 7B model never deadlocked, so counterfactual agents had no
effect. On tasks built to deadlock, they recover deadlocks and raise convergence success, at
about 2.8 times the model calls.

\section*{Individual member contributions}
{\small
\begin{tabularx}{\linewidth}{@{}l X@{}}
\toprule
Name & Contribution \\
\midrule
Naitik Gupta & Blackboard, PXP grammar and intelligibility; agent and provider pipeline; CI; Sections 1, 3. \\
Malay Neema & Evaluation design, metrics and statistics; benchmark ingest and scoring; Sections 4, 5. \\
Mayank Sharma & Scheduler, deadlock detection and counterfactual engine; compute setup and experiment runs; Section 6. \\
Yash Singh & Web interface; prompt library and personas; Section 2 and the video. \\
\bottomrule
\end{tabularx}}

\smallskip
\noindent All members have read this report and agree with the statement above.

\section*{Generative AI disclosure}
We used Claude Code (Anthropic, Claude Opus 5.5) as an assistant. Working from our
instructions, it drafted much of the code, tests and documentation (in particular for the
system rebuilt after the midsemester review), helped set up and run the experiments, and
produced a first draft of this report. The decisions were ours: the reading of the brief and
the PXP rules, the experimental design, the model and hardware, and the reduced scale.

We did not accept its output unchecked. Every change went through a pull request and was merged
only after the test suite (355 tests) and CI passed, and we rejected or redirected proposals that
did not fit the brief or our constraints. Our smoke runs on real data exposed defects in the
drafted code, each of which was traced and fixed with a test: a malformed MSCoRe file, a schema
that allowed empty judgments, example values that leaked into the prompts, and FHIR path
quoting. We checked every number in this report against the result files and revised the text,
and we take responsibility for all code and text. Qwen2.5-7B-Instruct is part of the method
(Section~3).

\section*{Code availability}
\noindent Code and data:\\
\url{\githuburl}\\
Setup is in the README and results are in \texttt{results/}. To reproduce, run \texttt{bb ablate} on \texttt{ablation-quick-ollama.toml} (language-model run) or \texttt{mock-synthetic.toml} (pipeline check) in \texttt{experiments/configs/}.

\begin{thebibliography}{99}
\small
\bibitem{erman1980}
L.~D. Erman, F.~Hayes-Roth, V.~R. Lesser and D.~R. Reddy.
\newblock The Hearsay-II speech-understanding system: integrating knowledge to resolve
uncertainty.
\newblock \emph{ACM Computing Surveys}, 12(2):213--253, 1980.

\bibitem{baskar2023}
A.~Baskar, A.~Srinivasan, M.~Bain and E.~Coiera.
\newblock A model for intelligible interaction between agents that predict and explain.
\newblock arXiv:2301.01819, 2023.

\bibitem{mestha2025}
H.~Mestha, K.~Bania, S.~V. Sathyanarayana, S.~Liu and A.~Srinivasan.
\newblock Multi-turn human--LLM interaction through the lens of a two-way intelligibility
protocol.
\newblock arXiv:2410.20600, 2025.

\bibitem{du2024debate}
Y.~Du, S.~Li, A.~Torralba, J.~B. Tenenbaum and I.~Mordatch.
\newblock Improving factuality and reasoning in language models through multiagent debate.
\newblock In \emph{ICML}, 2024.

\bibitem{foerster2018coma}
J.~Foerster, G.~Farquhar, T.~Afouras, N.~Nardelli and S.~Whiteson.
\newblock Counterfactual multi-agent policy gradients.
\newblock In \emph{AAAI}, 2018.

\bibitem{mscore}
Y.~Lei, H.~Xie, J.~Zhao, S.~Liu and X.~Song.
\newblock MSCoRe: a benchmark for multi-stage collaborative reasoning in LLM agents.
\newblock arXiv:2509.17628, 2025.

\bibitem{kramabench}
E.~Lai, G.~Vitagliano, Z.~Zhang et~al.
\newblock KramaBench: a benchmark for AI systems on data-to-insight pipelines over data lakes.
\newblock In \emph{ICLR}, 2026.

\bibitem{medagentbench}
Y.~Jiang, K.~C. Black, G.~Geng, D.~Park, J.~Zou, A.~Y. Ng and J.~H. Chen.
\newblock MedAgentBench: a realistic virtual EHR environment to benchmark medical LLM agents.
\newblock arXiv:2501.14654, 2025.

\bibitem{qwen25}
Qwen Team.
\newblock Qwen2.5 technical report.
\newblock arXiv:2412.15115, 2024.

\bibitem{kwon2023vllm}
W.~Kwon, Z.~Li, S.~Zhuang et~al.
\newblock Efficient memory management for large language model serving with PagedAttention.
\newblock In \emph{SOSP}, 2023.
\end{thebibliography}

\end{document}
